\chapter{A Concrete Naming Certificate for $\RegularCauchyTransformUp$}
\label{ch:concrete-instances-regular-cauchy-transform-namecert}
\origin{human}

Following Bishop and Bridges constructive analysis, the $\RegularCauchyTransformUp$ packet records the finite transform route between regular Cauchy names without committing to a map-specific analytic hypothesis. It sits over the regular rational readback of \autoref{ch:concrete-instances-regseqrat-namecert}, the finite observation windows of \autoref{ch:concrete-instances-streamname-namecert}, the dyadic tolerance ledger of \autoref{ch:concrete-instances-dyadicratcore-namecert}, the explicit modulus surface of \autoref{ch:concrete-instances-regular-modulus-namecert}, and the terminal real-name seal of \autoref{ch:concrete-instances-real-namecert}. The packet gives Real and Completion consumers a displayed transform certificate; it is not a quotient stream operation, selected representative, host function equality, countable-choice schedule, or ambient completeness principle.

\begin{definition}[Regular Cauchy transform carrier]
\label{def:regular-cauchy-transform-carrier}
A $\RegularCauchyTransformUp$ carrier is a finite $\BHist$ packet
$$
T=(S,R,M,D,U,Q,E,H,C,P,N)
$$
with the following displayed rows:
\begin{enumerate}
\item $S$ is the source $\StreamNameUp$ finite-window row.
\item $R$ is the source $\RegSeqRatUp$ regular Cauchy readback row attached to those windows.
\item $M$ is the transform modulus row, assigning each requested output tolerance to a finite source-window obligation.
\item $D$ is the $\DyadicRatCoreUp$ tolerance ledger carrying the requested output radius, source radius, and slack comparisons used by $M$.
\item $U$ is the output $\StreamNameUp$ finite-window row produced by the displayed transform route.
\item $Q$ is the output $\RegSeqRatUp$ readback row attached to $U$ and $D$.
\item $E$ is the terminal $\RealUp$ seal reached only after the output regular readback row has been displayed.
\item $H$ records componentwise $\hsame$ transport for the source windows, source readback, transform modulus, dyadic ledger, output windows, output readback, Real seal, replay, provenance, and local naming rows.
\item $C$ records displayed $\Cont$ replay of the route
$$
S,R,M,D \longmapsto U,Q \longmapsto E .
$$
\item $P$ is the $\Pkg$ provenance over $\RegularCauchyTransformUp$, $\StreamNameUp$, $\RegSeqRatUp$, $\DyadicRatCoreUp$, $\RegularModulusUp$, $\RealUp$, $\BHist$, $\hsame$, $\Cont$, and $\NameCert$.
\item $N$ is the local $\NameCert_{\RegularCauchyTransformUp}$ row.
\end{enumerate}
The carrier has no coordinate for quotient stream equality, a selected Cauchy representative, host function extensionality, countable choice, ambient $\RealUp$ completeness, or a detached output seal.
\end{definition}

\begin{definition}[Regular Cauchy transform classifier]
\label{def:regular-cauchy-transform-classifier}
The $\RegularCauchyTransformUp$ classifier compares accepted carriers only through componentwise $\hsame$ transport of
$$
S,R,M,D,U,Q,E,H,C,P,N.
$$
Thus source windows, source regular readback, transform modulus, dyadic tolerance accounting, output windows, output regular readback, terminal Real seal, replay, provenance, and local naming are all visible comparison data. No classifier clause compares hidden host functions, quotient streams, selected limits, or unrecorded completion witnesses.
\end{definition}

\begin{theorem}[Regular Cauchy transform NameCert obligations]
\label{thm:regular-cauchy-transform-namecert-obligations}
For every accepted $\RegularCauchyTransformUp$ carrier
$$
T=(S,R,M,D,U,Q,E,H,C,P,N),
$$
the local $\NameCert_{\RegularCauchyTransformUp}$ surface consists exactly of carrier admission, source-window exposure, source regular readback, transform-modulus routing, dyadic tolerance accounting, output-window production, output regular readback, terminal Real-seal handoff, componentwise transport, displayed replay, provenance, and local naming through the displayed rows $S,R,M,D,U,Q,E,H,C,P,N$.
\end{theorem}
\begin{proof}
Project the carrier of \autoref{def:regular-cauchy-transform-carrier}. The source side is visible as $S,R$, the tolerance route is visible as $M,D$, the output side is visible as $U,Q,E$, and the structural rows are visible as $H,C,P,N$.

Read the transform route in displayed order. A consumer first opens finite source windows, attaches them to the regular rational readback, reads the transform modulus and dyadic tolerance ledger, and only then obtains output windows and output regular readback.

Finally attach the terminal and structural rows. The Real seal $E$ is downstream of $Q$, while $H,C,P,N$ transport, replay, source, and name the same finite route. Because the carrier has no hidden host map, quotient stream equality, selected representative, choice schedule, or ambient completeness row, the listed obligations exhaust the local NameCert surface.
\end{proof}

\begin{theorem}[Regular Cauchy transform modulus route]
\label{thm:regular-cauchy-transform-modulus-route}
Every output window admitted by a $\RegularCauchyTransformUp$ packet is backed by the transform modulus row $M$, the dyadic tolerance ledger $D$, a finite source $\StreamNameUp$ window $S$, and source $\RegSeqRatUp$ readback $R$. Hence a Real or Completion consumer cannot read the output regular name through $U,Q$ without first displaying the modulus-controlled source route.
\end{theorem}
\begin{proof}
Begin with an admitted output window. Carrier projection exposes $U$ as an output $\StreamNameUp$ row, but $U$ is downstream of the transform route rather than an independent stream source.

Trace the request through $M$ and $D$. The modulus row names the source-window obligation for the requested output tolerance, and the dyadic ledger records the radii and slack comparisons needed to serve that obligation.

The source evidence is then $S,R$. The finite windows in $S$ are attached to regular rational readback by $R$, so the output route cannot be justified by a selected limit or host convergence witness.

After $U$ is produced, the output readback $Q$ and terminal seal $E$ remain downstream of the same route. The structural rows $H,C,P,N$ preserve this order under transport and replay.
\end{proof}

\begin{theorem}[Regular Cauchy transform Real-seal nonescape]
\label{thm:regular-cauchy-transform-real-seal-nonescape}
Every terminal $\RealUp$ read accepted by a $\RegularCauchyTransformUp$ carrier factors through
$$
\StreamNameUp,\ \RegSeqRatUp,\ \RegularModulusUp,\ \DyadicRatCoreUp,\ \StreamNameUp,\ \RegSeqRatUp,\ \RealUp,\ \BHist,\ \hsame,\ \Cont,\ \Pkg,\ \NameCert .
$$
No accepted read bypasses the source regular name, transform modulus, dyadic tolerance ledger, output regular readback, componentwise transport, replay, provenance, or local naming rows.
\end{theorem}
\begin{proof}
Start with the terminal seal row $E$. In the carrier definition, $E$ is reachable only after the output readback row $Q$, and $Q$ is attached to the output finite-window row $U$.

Use the modulus-route theorem. The output window $U$ is produced only after the source windows $S$, source readback $R$, transform modulus $M$, and dyadic ledger $D$ have been displayed.

Finally read the structural rows. The $\hsame$ transport, $\Cont$ replay, $\Pkg$ provenance, and local $\NameCert$ preserve the same finite chain; they do not add a quotient stream, selected host representative, ambient completeness theorem, or second seal route.
\end{proof}

\falsifiablePrediction{A $\RegularCauchyTransformUp$ consumer that reaches a terminal $\RealUp$ seal without first displaying source windows, source RegSeqRat readback, a transform modulus row, dyadic tolerance accounting, output windows, output RegSeqRat readback, transport, replay, provenance, and local naming is not reading this packet.}

\independenceWitness{regseqrat, streamname, dyadicratcore, regular-modulus, real: no listed sibling supplies the whole transform packet because each contributes only readback, windowing, tolerance, modulus, or sealing data while this chapter binds the finite source-to-output route.}

\closureat{RegularCauchyTransformUp}{seedStr}

\begin{closurestatus}{\RegularCauchyTransformUp}
  \constructivestory{A $\RegularCauchyTransformUp$ packet is generated from source StreamName windows, source RegSeqRat readback, a transform modulus row, DyadicRatCore tolerance accounting, output StreamName windows, output RegSeqRat readback, a terminal RealUp seal, componentwise hsame transport, displayed Cont replay, Pkg provenance, and a local NameCert row.}
  \theoryclosure{\seedClosure}
  \scopeclosed{\autoref{def:regular-cauchy-transform-carrier}, \autoref{def:regular-cauchy-transform-classifier}, \autoref{thm:regular-cauchy-transform-namecert-obligations}, \autoref{thm:regular-cauchy-transform-modulus-route}, and \autoref{thm:regular-cauchy-transform-real-seal-nonescape} seed the finite transform route between regular Cauchy names.}
  \formalstatus{\unformalizedV}
  \bridgestatus{none}
  \notclaimed{This chapter does not close quotient stream equality, selected representatives, host function equality, countable choice, ambient Real completeness, Lean verification, or bridge checking.}
  \upgradepath{A formal route can encode the displayed carrier and prove carrier admission, classifier transport, modulus routing, output-window production, Real-seal handoff, and ledger non-escape for BEDC.Derived.RegularCauchyTransformUp.}
\end{closurestatus}
